\section{Overview}

This chapter presents the MEVRES mechanism. We first explain how it departs from RediSwap, then define the model, describe the five protocol steps, and prove the properties claimed for the mechanism.

\subsection{Motivation: the Sequencer in RediSwap}

RediSwap~\cite{zhang2024rediswap} sells MEV after the fact and refunds the proceeds. It receives transactions through a private channel, which creates two problems.

\paragraph{A trusted intermediary.} Whoever operates the channel sees every transaction in plaintext before the auction.

\paragraph{A single point of failure.} An operator who is malicious or compromised can extract the MEV directly, reveal transactions to a favoured arbitrageur, or otherwise bias the auction. Users and liquidity providers then receive no refund.

\paragraph{Our approach.} MEVRES replaces the sequencer with a commitment and a threshold encryption layer. A transaction stays unreadable to every party, including the arbitrageurs and the contract that runs the auction, until all bids are fixed. Decryption requires the cooperation of $t+1$ out of $n$ independent nodes, so no single party can open a transaction early. The auction and the refund are RediSwap's; the contribution of MEVRES is to make them hold without a trusted party.

Figure~\ref{fig:protocol-flow} gives an overview of the protocol.

\begin{figure}[h]
\centering
\includegraphics[width=\textwidth]{protocol_flow.pdf}
\caption{The five steps of MEVRES. The transaction is hidden during Steps 1 and 2 and is opened only after bids are fixed. Steps 3 to 5 are enforced by the smart contract.}
\label{fig:protocol-flow}
\end{figure}

\section{Formal Model}

\subsection{Participants}

The mechanism involves five classes of participants.
\begin{itemize}
    \item \textbf{Users} (noise traders), who submit swap transactions.
    \item \textbf{Arbitrageurs} (informed traders), who bid for MEV opportunities. Arbitrageur $i\in[m]$ holds a private belief $v_i^{*}$ about the external price of $X$.
    \item \textbf{Liquidity providers}, who own the pool reserves and receive swap fees and part of the auction revenue.
    \item \textbf{Threshold nodes} $N_1,\dots,N_n$, which jointly hold the decryption key.
    \item \textbf{The smart contract}, which records submissions, runs the auction, executes bundles and pays refunds.
\end{itemize}

\subsection{Trust and Threat Model}

The adversary may observe all on-chain data and all network traffic, may control any number of arbitrageurs and the block producer, and may corrupt up to $t$ of the $n$ threshold nodes. We assume:
\begin{itemize}
    \item $H$ is a collision-resistant hash function, modelled as a random oracle;
    \item the threshold encryption scheme is semantically secure against any coalition of at most $t$ nodes~\cite{desmedt1989threshold};
    \item the threshold key is generated by a distributed key generation protocol, so that no party ever holds the full secret key~\cite{pedersen1991threshold};
    \item at least $t+1$ nodes are honest and online in every block, which requires $n\geq 2t+1$;
    \item the smart contract executes correctly and atomically.
\end{itemize}
Privacy therefore holds as long as at most $t$ nodes are corrupt, and liveness holds as long as at least $t+1$ nodes respond.

\subsection{Valuation of a Sandwich Bundle}

Let $s_0=(x_0,y_0)$ be the pool state at the start of the block. Consider a bundle built around a user transaction $TX_j$ that starts and ends at $s_0$. Because the pool's reserves are unchanged by the bundle, the arbitrageur who executes it ends up holding exactly what the user gave up and owing exactly what the user received. For an arbitrageur who values $X$ at price $v$, the profit of the bundle is therefore
\begin{equation}
\Pi(v,\,TX_j,\,s_0)=
\begin{cases}
v\cdot\delta^{x}_{in}-\delta^{y}_{out}, & TX_j=(X\rightarrow Y,\ \delta^{x}_{in},\ \delta^{y}_{out}),\\[4pt]
\delta^{y}_{in}-v\cdot\delta^{x}_{out}, & TX_j=(Y\rightarrow X,\ \delta^{y}_{in},\ \delta^{x}_{out}),
\end{cases}
\label{eq:pi}
\end{equation}
measured in units of $Y$. The true value of the opportunity to arbitrageur $i$ is
\begin{equation}
v_i(TX_j)=\Pi(v_i^{*},\,TX_j,\,s_0).
\label{eq:value}
\end{equation}
Note that $\Pi$ is strictly increasing in $v$ for an $X\rightarrow Y$ trade and strictly decreasing in $v$ for a $Y\rightarrow X$ trade.

\subsection{Auction Design}

Each arbitrageur submits a single price report $q_i$ per block. The arbitrageur does not need to know $TX_j$: once the transaction is decrypted, the contract converts every report into a bid
\begin{equation}
b_{i,j}=\Pi(q_i,\,TX_j,\,s_0).
\label{eq:bid}
\end{equation}
The winner of the opportunity is the arbitrageur with the highest bid, and the winner pays the second-highest bid:
\begin{equation}
w_j=\arg\max_{i\in[m]} b_{i,j},
\qquad
p_j=\max\Bigl(0,\ \max_{i\neq w_j} b_{i,j}\Bigr).
\label{eq:winner}
\end{equation}
If no bid is positive, no bundle is formed and the transaction is not executed in this block. Ranking by $b_{i,j}$ rather than by $q_i$ matters: for an $X\rightarrow Y$ trade the highest report gives the highest bid, whereas for a $Y\rightarrow X$ trade the lowest report does.

\section{The Five-Step Protocol}

All five steps complete within one block.

\subsection{Step 1: Commit and Encrypt}

The user fixes the trade with a hash commitment and hides it with threshold encryption. Let $pk$ be the public key of the threshold scheme.

\begin{algorithm}[h]
\caption{User Submission (Step 1)}
\begin{algorithmic}[1]
\Require Trade $TX=(d,\ \delta_{in},\ \delta_{out})$ with direction $d\in\{X\rightarrow Y,\ Y\rightarrow X\}$; threshold public key $pk$
\Ensure On-chain record $(c,\ CT)$
\State User samples a random salt $r\xleftarrow{\$}\{0,1\}^{256}$
\State User computes the commitment $c\leftarrow H(TX\,\|\,r)$
\State User computes the ciphertext $CT\leftarrow \mathsf{Enc}_{pk}(TX\,\|\,r)$
\State User submits $(c,\ CT)$ together with collateral $D$
\State Contract records $(c,\ CT)$ and emits \texttt{TxSubmitted}
\end{algorithmic}
\end{algorithm}

The commitment is to the \emph{plaintext} transaction, and the salt is encrypted together with it. This has three consequences. First, the salt gives the committed value 256 bits of entropy, so an observer cannot test guesses of $TX$ against $c$. Second, the user does not have to come back and reveal anything: decryption in Step 3 recovers both $TX$ and $r$. Third, the contract can check that what the nodes decrypt is what the user committed to. The collateral $D$ is returned when this check passes and forfeited when it fails, which removes any incentive to submit ciphertexts that do not open correctly.

\subsection{Step 2: Blind Bidding}

While every transaction in the block is still encrypted, the arbitrageurs report their price beliefs.

\begin{algorithm}[h]
\caption{Arbitrageur Bidding (Step 2)}
\begin{algorithmic}[1]
\Require Belief $v_i^{*}$ of arbitrageur $i$; initial state $s_0$; ciphertexts $\{CT_j\}$
\Ensure On-chain report $q_i$
\State Arbitrageur $i$ observes $s_0$ and the number of pending ciphertexts
\State Arbitrageur $i$ submits a price report $q_i$ before the bid deadline
\State Contract records $q_i$; reports cannot be changed after the deadline
\end{algorithmic}
\end{algorithm}

An arbitrageur reports a \emph{price}, not an amount. This is what allows bidding on transactions it cannot see: the report is turned into a bid by Equation~\ref{eq:bid} only after decryption.

\subsection{Step 3: Decrypt, Verify and Select the Winner}

After the bid deadline each threshold node publishes a decryption share for every ciphertext. Any $t+1$ valid shares suffice to recover the plaintext.

\begin{algorithm}[h]
\caption{Decryption and Winner Selection (Step 3)}
\begin{algorithmic}[1]
\Require $(c_j,\ CT_j)$; reports $\{q_i\}$; decryption shares from the threshold nodes
\State Contract combines $t+1$ decryption shares to obtain $(TX_j,\ r_j)$
\If{fewer than $t+1$ valid shares arrive by the deadline}
    \State $TX_j$ is cancelled, the collateral is returned and nothing is revealed
\ElsIf{$H(TX_j\,\|\,r_j)\neq c_j$}
    \State $TX_j$ is rejected and the user's collateral is forfeited
\Else
    \State Contract computes $b_{i,j}\leftarrow\Pi(q_i,\ TX_j,\ s_0)$ for every arbitrageur $i$
    \State Contract sets $w_j$ and $p_j$ according to Equation~\ref{eq:winner}
\EndIf
\end{algorithmic}
\end{algorithm}

\subsection{Step 4: Execute the Sandwich Bundle}

For each accepted transaction the contract builds an atomic bundle for the winner $w_j$:
\begin{equation}
\mathcal{B}_j=\langle \mathit{FrontRun}_j,\ TX_j,\ \mathit{BackRun}_j\rangle .
\label{eq:bundle}
\end{equation}
\begin{itemize}
    \item $\mathit{FrontRun}_j$: $w_j$ trades to move the pool from $s_0$ to the state $s_1$ from which the user's swap returns exactly $\delta_{out}$.
    \item $TX_j$: the user's swap executes, moving the pool from $s_1$ to $s_2$. The user receives the stated minimum output.
    \item $\mathit{BackRun}_j$: $w_j$ trades to move the pool from $s_2$ back to $s_0$.
\end{itemize}
For the constant product pool and an $X\rightarrow Y$ trade, the front-run state has reserve
\begin{equation}
x_1=\frac{-\delta^{x}_{in}+\sqrt{(\delta^{x}_{in})^{2}+4k\,\delta^{x}_{in}/\delta^{y}_{out}}}{2},
\qquad y_1=\frac{k}{x_1},
\label{eq:frontstate}
\end{equation}
which follows from requiring $y_1-k/(x_1+\delta^{x}_{in})=\delta^{y}_{out}$. Figure~\ref{fig:sandwich-curve} shows the three legs on the curve.

\begin{figure}[h]
\centering
\includegraphics[width=0.68\textwidth]{sandwich_curve.pdf}
\caption{The sandwich bundle on a constant product curve for an $X\rightarrow Y$ trade. The trade size is exaggerated for visibility. The bundle starts and ends at $s_0$.}
\label{fig:sandwich-curve}
\end{figure}

The bundle is a sandwich by construction. MEVRES does not stop the sandwich from happening; it decides who may perform it, guarantees that the user still receives their stated minimum, and makes the performer pay for the privilege.

\subsection{Step 5: Redistribute}

The payment $p_j$ is split between the user and the liquidity providers:
\begin{equation}
\text{Refund to user}_j=\alpha\cdot p_j,
\qquad
\text{Refund to LPs}=(1-\alpha)\cdot p_j,
\label{eq:refund}
\end{equation}
where $\alpha\in[0,1]$ is a protocol parameter. For sandwich MEV, $\alpha$ is set close to 1, because the value originates in the user's slippage tolerance. The swap fee on the user's input is paid to LPs as usual. The refund is paid inside the same bundle, so it cannot be intercepted.

\paragraph{Rebalancing MEV.} After all user bundles the pool is still at $s_0$, which may differ from the external price. The opportunity to move the pool to the external price is auctioned with the same reports: the bid of arbitrageur $i$ is its profit from moving the pool to price $q_i$, the highest bid wins and pays the second-highest bid, and the whole payment goes to LPs. This compensates LPs for loss-versus-rebalancing~\cite{milionis2022lvr}.

\section{Properties of the Mechanism}
\label{sec:properties}

\begin{theorem}[Bundle independence]
\label{thm:independence}
Let $\mathcal{B}_j$ and $\mathcal{B}_k$ be the bundles of two distinct user transactions in the same block. The value of each bundle to any arbitrageur does not depend on whether, or in which order, the other bundle is executed.
\end{theorem}
\begin{proof}
Every bundle starts at $s_0$ and, by construction of the back-run, ends at $s_0$. Whichever bundle executes second therefore also starts from $s_0$. By Equation~\ref{eq:pi}, the value of a bundle depends only on the arbitrageur's price, the transaction and the fact that the bundle starts and ends at the same state. It does not depend on any other transaction.
\end{proof}

\begin{theorem}[Incentive compatibility]
\label{thm:ic}
For every arbitrageur, reporting $q_i=v_i^{*}$ is a dominant strategy.
\end{theorem}
\begin{proof}
Fix a transaction $TX_j$ and the reports of all other arbitrageurs, and let $\beta=\max_{k\neq i}b_{k,j}$. If arbitrageur $i$ wins, its utility is $v_i(TX_j)-\max(0,\beta)$, which does not depend on its own report. If it loses, its utility is zero. A truthful report gives the bid $b_{i,j}=v_i(TX_j)$, so $i$ wins exactly when $v_i(TX_j)>\max(0,\beta)$, that is, exactly when winning is profitable. Any other report either leaves the outcome unchanged, or wins an opportunity whose price exceeds its value, or loses one whose value exceeds its price. This is the standard argument for the second-price auction~\cite{vickrey1961counterspeculation}. By Theorem~\ref{thm:independence} the opportunities in a block are independent, so the single report $q_i=v_i^{*}$ is optimal for all of them simultaneously.
\end{proof}

\begin{theorem}[Individual rationality]
\label{thm:ir}
No participant is made worse off by taking part.
\begin{enumerate}
    \item A truthful arbitrageur never has negative utility.
    \item A user whose transaction executes receives at least the minimum output they specified.
    \item The pool reserves are unchanged by a bundle, and liquidity providers receive a non-negative payment.
\end{enumerate}
\end{theorem}
\begin{proof}
(1) A truthful arbitrageur wins only when $v_i(TX_j)>p_j$, and has utility zero otherwise. (2) The front-run state is chosen so that the user's swap returns exactly $\delta_{out}$, and the refund $\alpha p_j$ is non-negative because $p_j\geq 0$. (3) The bundle returns the pool to $s_0$, and LPs receive the swap fee and $(1-\alpha)p_j\geq 0$.
\end{proof}

\begin{theorem}[Bidder-side Sybil-proofness]
\label{thm:sybil}
An arbitrageur cannot increase its utility by submitting reports under several identities.
\end{theorem}
\begin{proof}
Let $\beta$ be the highest bid among the other arbitrageurs. Whichever identity wins, the price it pays is the second-highest bid among \emph{all} bids, which is at least $\max(0,\beta)$. The arbitrageur's utility is therefore at most $v_i(TX_j)-\max(0,\beta)$ when one of its identities wins, and zero otherwise. A single truthful report achieves $\max\bigl(0,\ v_i(TX_j)-\max(0,\beta)\bigr)$, which is at least as large.
\end{proof}

\begin{proposition}[Pre-auction privacy]
\label{prop:privacy}
Under the assumptions of the threat model, no coalition of arbitrageurs, block producers and at most $t$ threshold nodes learns anything about a transaction, other than the length of its ciphertext, before the bid deadline.
\end{proposition}
\begin{proof}[Proof sketch]
Before the deadline the coalition sees only $c$ and $CT$. The ciphertext reveals nothing by the semantic security of the threshold scheme against $t$ corrupt nodes. The commitment $c=H(TX\,\|\,r)$ reveals nothing because $r$ is uniformly random, unknown to the coalition, and $H$ is a random oracle. A full proof that treats the two primitives together under an adaptive adversary is left to future work.
\end{proof}

Theorem~\ref{thm:sybil} concerns arbitrageurs only. We do not claim that a user cannot benefit from splitting one trade into several, and we do not analyse collusion among arbitrageurs; both are discussed in Chapter~\ref{ch:conclude}.

\section{Comparison with RediSwap}

Table~\ref{tab:rediswap} summarises the differences between the two mechanisms.

\begin{table}[h]
\centering
\caption{Detailed Comparison: RediSwap vs. MEVRES}
\label{tab:rediswap}
\begin{tabular}{|L{3.9cm}|L{4.3cm}|L{4.7cm}|}
\hline
\textbf{Property} & \textbf{RediSwap} & \textbf{MEVRES} \\
\hline
Transaction privacy & Private channel run by a sequencer & Hash commitment and threshold encryption \\
\hline
Trust assumption & The sequencer is honest & At most $t$ of $n$ nodes are corrupt \\
\hline
Who cannot front-run & External searchers & All parties, including the sequencer role and the block producer \\
\hline
MEV redistribution & Yes (users and LPs) & Yes (users and LPs) \\
\hline
Incentive compatibility & Yes & Yes (Theorem~\ref{thm:ic}) \\
\hline
Single point of failure & Yes (the sequencer) & No \\
\hline
Liveness requirement & The sequencer is online & $t+1$ nodes are online \\
\hline
\end{tabular}
\end{table}
